% TOP_PROBLEM: {"id": 2566, "title": "Kourovka Notebook Problem 21.57", "category_id": 17, "category": {"id": 17, "name": "group_theory", "display_name": "Group Theory", "description": "Problems about groups, group actions, representations, and related algebraic structures.", "slug": "group-theory", "order_index": 17, "created_at": "2026-07-31T15:26:25.671Z"}, "status": "open", "problem_number": "KOU-21.57", "set_id": 10}
% FIRST_POSTED: 2026-10-01
% ENTRY_KIND: statement_only
% UnsolvedMath ID 2566; source catalogue code KOU-21.57
% !TeX program = lualatex
% Definition and sources only; this file is not an LLM solution attempt.
\documentclass[11pt,a4paper]{article}
\usepackage[margin=25mm]{geometry}
\usepackage{fontspec}
\setmainfont{lmroman10-regular.otf}[BoldFont=lmroman10-bold.otf,
 ItalicFont=lmroman10-italic.otf,BoldItalicFont=lmroman10-bolditalic.otf]
\usepackage{amsmath,amssymb,mathtools}
\usepackage{unicode-math}
\setmathfont{latinmodern-math.otf}
\AtBeginDocument{\let\setminus\smallsetminus}
\usepackage{xurl,hyperref}
\hypersetup{colorlinks=true,urlcolor=blue}
\setcounter{secnumdepth}{0}
\setlength{\parindent}{0pt}
\setlength{\parskip}{5pt}
\setlength{\emergencystretch}{3em}
\begin{document}
\section*{Kourovka Notebook Problem 21.57}\label{2566}
{\small UnsolvedMath ID 2566 · KOU-21.57 · Group Theory · Kourovka Notebook - New Problems, Issue 21}
\subsection{Definitions and mathematical statement}
Let $\mathcal X$ be a nonempty class of finite groups, each of odd
order. Assume it is closed under isomorphism, subgroups, quotients and
extensions. The last condition means that if $N\trianglelefteq K$ and
both $N$ and $K/N$ belong to $\mathcal X$, then $K$ belongs to
$\mathcal X$. Because the class is nonempty and subgroup-closed,
it contains the trivial group.

For a finite group $G$, an $\mathcal X$-maximal subgroup is a subgroup
$H\le G$ such that $H\in\mathcal X$ and there is no strictly larger
subgroup $J\le G$ belonging to $\mathcal X$.
This is maximality within the specified class, not necessarily maximality
among all proper subgroups. If $G\in\mathcal X$, the subgroup $G$
itself is $\mathcal X$-maximal.

Given an $\mathcal X$-maximal subgroup $H$ of any finite group $G$
and a normal subgroup $N\trianglelefteq G$, must $H\cap N$ be
$\mathcal X$-maximal in $N$? In explicit terms, does
\[
 H\cap N\le J\le N,\quad J\in\mathcal X
 \quad\Longrightarrow\quad J=H\cap N
\]
always hold?

Membership $H\cap N\in\mathcal X$ already follows from subgroup
closure. The issue is its maximality. The ambient finite group may
have even order and need not belong to $\mathcal X$. The odd-order
restriction is imposed on the class members, and all allowed classes,
normal subgroups and class-maximal choices of $H$ are quantified over.
\subsection{Short English statement}
Does intersecting a class-maximal subgroup with a normal subgroup preserve maximality for every extension-closed class of finite odd-order groups?
\subsection{Sources}
\begin{itemize}
\item[S1] ULAM.AI and the UnsolvedMath Contributors, supplied problems.json, record 2566 (KOU-21.57), Kourovka Notebook - New Problems, Issue 21. Catalogue identity and source transcription; CC BY 4.0.
\url{https://huggingface.co/datasets/ulamai/UnsolvedMath}.
\item[S2] The Kourovka Notebook, 21st edition, new problems, Problem 21.57. Expanded formulation of the supplied catalogue statement.
\url{https://arxiv.org/pdf/1401.0300v47}.
\end{itemize}
\end{document}
